% ============================================================
% The complement of Q
% Companion note to monograph_v2.tex, Part III
% Author: R. Wydaeghe
% ============================================================

\documentclass[11pt,a4paper]{article}

\usepackage[margin=2.4cm]{geometry}
\usepackage[utf8]{inputenc}
\usepackage[T1]{fontenc}
\usepackage{lmodern}
\usepackage{microtype}
\linespread{1.04}

\usepackage{amsmath,amssymb,amsthm,mathtools}
\usepackage{bm}
\usepackage{mathrsfs}
\usepackage{bbm}
\usepackage{cancel}
\usepackage{xfrac}
\allowdisplaybreaks[2]

\usepackage{empheq}
\usepackage[most]{tcolorbox}
\tcbset{highlight math style={%
    enhanced, colframe=black!70, colback=gray!6,
    arc=2pt, boxrule=0.5pt, left=3pt, right=3pt}}

\usepackage{siunitx}
\sisetup{detect-all, range-phrase = {--}, range-units = single,
         per-mode = symbol, inter-unit-product = {\,}}

\usepackage{enumitem}
\usepackage{xcolor}
\usepackage{xspace}
\usepackage{fancyhdr}
\pagestyle{fancy}
\fancyhf{}
\cfoot{\thepage}
\renewcommand{\headrulewidth}{0pt}

\usepackage{titlesec}
\titleformat{\section}{\large\bfseries}{\thesection}{1em}{}
\titleformat{\subsection}{\normalsize\bfseries}{\thesubsection}{1em}{}

\usepackage[colorlinks=true, linkcolor={blue!55!black},
            citecolor={green!45!black}, urlcolor={blue!65!black}]{hyperref}
\usepackage[nameinlink,capitalize]{cleveref}

\usepackage{aliascnt}
\theoremstyle{definition}
\newtheorem{theorem}{Theorem}[section]
\newaliascnt{proposition}{theorem}
\newtheorem{proposition}[proposition]{Proposition}
\aliascntresetthe{proposition}
\newaliascnt{lemma}{theorem}
\newtheorem{lemma}[lemma]{Lemma}
\aliascntresetthe{lemma}
\newaliascnt{corollary}{theorem}
\newtheorem{corollary}[corollary]{Corollary}
\aliascntresetthe{corollary}
\newaliascnt{definition}{theorem}
\newtheorem{definition}[definition]{Definition}
\aliascntresetthe{definition}
\newaliascnt{remark}{theorem}
\newtheorem{remark}[remark]{Remark}
\aliascntresetthe{remark}
\crefname{proposition}{Proposition}{Propositions}
\Crefname{proposition}{Proposition}{Propositions}
\crefname{lemma}{Lemma}{Lemmas}\Crefname{lemma}{Lemma}{Lemmas}
\crefname{corollary}{Corollary}{Corollaries}\Crefname{corollary}{Corollary}{Corollaries}
\crefname{definition}{Definition}{Definitions}\Crefname{definition}{Definition}{Definitions}
\crefname{remark}{Remark}{Remarks}\Crefname{remark}{Remark}{Remarks}

% Custom macros (aligned with monograph)
\newcommand{\khat}{\hat{\bm{k}}}
\newcommand{\nhat}{\hat{\bm{n}}}
\newcommand{\rr}{\mathbf{r}}
\newcommand{\EE}{\mathbf{E}}
\newcommand{\Sinc}{S_{\mathrm{inc}}}
\newcommand{\Sab}{S_{\mathrm{ab}}}
\newcommand{\Sref}{S_{\mathrm{re}}}
\newcommand{\Sinn}{S_{\mathrm{in}}}
\newcommand{\Smis}{S_{\mathrm{mi}}}
\DeclareMathOperator{\ReLU}{ReLU}
\DeclareMathOperator{\RE}{Re}
\DeclareMathOperator{\IM}{Im}
\DeclareMathOperator{\tr}{tr}
\newcommand{\diff}{\mathrm{d}}
\newcommand{\ident}{\mathbbm{1}}
\newcommand{\Reals}{\mathbb{R}}
\newcommand{\Complex}{\mathbb{C}}
\newcommand{\half}{\tfrac{1}{2}}
\newcommand{\Gtil}{\tilde{\mathbf{G}}}
\newcommand{\gtil}{\tilde{\mathbf{g}}}
\newcommand{\Fn}{\mathbf{F}_{\!n}}
\newcommand{\Rn}{\mathbf{R}_{\!n}}
\newcommand{\Qabs}{\mathbf{Q}_{\mathrm{ab}}}
\newcommand{\Qref}{\mathbf{Q}_{\mathrm{re}}}
\newcommand{\Qinn}{\mathbf{Q}_{\mathrm{in}}}
\newcommand{\Qmis}{\mathbf{Q}_{\mathrm{mi}}}

\title{\textbf{The complement of $\mathbf{Q}$:\\
scattered power, transparent precoders, thermal duality}}
\author{Robin Wydaeghe\\
 Ghent University \& imec\\
 \texttt{robin.wydaeghe@ugent.be}}
\date{\today}

\makeatletter
\renewcommand{\maketitle}{%
  \thispagestyle{plain}%
  \begin{center}
    \rule{\linewidth}{1pt}\\[0.4em]
    {\LARGE\bfseries \@title}\\[0.5em]
    {\large \@author}\\[0.3em]
    {\small \@date}\\[0.4em]\rule{\linewidth}{1pt}
  \end{center}
  \vspace{0.5em}
}
\makeatother

\begin{document}
\maketitle

\begin{abstract}
In Part III of the monograph the whole-body absorbed power is
$P_{\mathrm{ab}}=\mathbf{x}^H\Qabs\mathbf{x}$, with $\Qabs$ the
Fresnel-transmission Gram integral of the exposure channel.
Power conservation asks a question this operator does not answer:
of the $\|\mathbf{x}\|^2$ watts that leave the array, what fraction
is \emph{not} absorbed, and where does it go?
This note writes the book-keeping explicitly.  We introduce three
companion Hermitian operators $\Qref$, $\Qinn$, $\Qmis$, show that
$\Qinn=\Qabs+\Qref$ is an exact consequence of Fresnel energy
conservation, and that
$\Qabs+\Qref+\Qmis\preceq\mathbf{I}$ is the array-level version of
``what goes out must come back, be absorbed, or leave forever''.
Under the pseudo-Brewster regime $\Qabs$ and $\Qref$ share
eigenvectors, so every exposure-constrained precoder is
simultaneously a scattering-optimal waveform: the worst beam for
the bystander is the best beam for a monostatic radar.  We finish
by sketching the thermal-reciprocal reading of $\Qabs$ (the
Kirchhoff dual) and the generic existence of a \emph{transparency
subspace} $\ker\Qinn$ whenever $M>3M_{\triangle}$.
\end{abstract}

\vspace{0.3em}
\noindent\textit{Style note.}~This is a wandering companion to
Part~III, not a theorem-proof paper.  Sections are short.  Some
claims are exact, others approximate in the pseudo-Brewster sense,
others are conjectures flagged as such.

\tableofcontents
\bigskip

% ============================================================
\section{Where does the rest of the power go?}
% ============================================================

From the exposure-operator definition in Part~III of the monograph,
\begin{equation}
  P_{\mathrm{ab}} = \mathbf{x}^H\Qabs\mathbf{x},
  \qquad
  \Qabs = \int_\Sigma \Gtil(\rr)^H\Gtil(\rr)\,\diff A.
\end{equation}
The array radiates a total power that, in the convention of the
monograph where $\mathbf{x}\in\Complex^M$ is scaled so that each
entry is an amplitude at the element feed, is
\begin{equation}\label{eq:Ptx}
  P_{\mathrm{tx}} = \|\mathbf{x}\|^2 = \mathbf{x}^H\mathbf{I}\mathbf{x}.
\end{equation}
The ratio $\eta_{\mathrm{ab}}(\mathbf{x})
=\mathbf{x}^H\Qabs\mathbf{x}/\|\mathbf{x}\|^2\in[0,1]$ is the
whole-body absorption efficiency of the precoder.  What is
$1-\eta_{\mathrm{ab}}$?

There are three places the un-absorbed watts can go:
\begin{enumerate}[leftmargin=1.6em,itemsep=2pt]
\item they hit the body and bounce: Fresnel reflection back into
      the environment;
\item they hit the body and ride along the surface: evanescent
      tangential fields and Norton-wave contributions that we drop
      here, see \cref{rem:ignored};
\item they never hit the body at all: radiation that lights up the
      sky, the ceiling, the neighbouring building, or the inside of
      a reverberation chamber.
\end{enumerate}
We want a matrix for each.  Three operators, one identity.

% ============================================================
\section{The reflection operator}
% ============================================================

\subsection{Fresnel amplitude reflection, revived}

The monograph uses amplitude transmission coefficients
\begin{equation}\label{eq:t}
  t_{s,n}=\frac{2\mu_n}{\mu_n+\xi_n},\qquad
  t_{p,n}=\frac{2\tilde n\,\mu_n}{\tilde n^2\mu_n+\xi_n},
\end{equation}
with $\xi_n=\sqrt{\tilde n^2-1+\mu_n^2}$ and $\RE\xi_n>0$.
The \emph{reflection} siblings are
\begin{equation}\label{eq:r}
  r_{s,n}=\frac{\mu_n-\xi_n}{\mu_n+\xi_n},\qquad
  r_{p,n}=\frac{\tilde n^2\mu_n-\xi_n}{\tilde n^2\mu_n+\xi_n},
\end{equation}
linked to \cref{eq:t} by $t_{s,n}=1+r_{s,n}$ and
$t_{p,n}=(1+r_{p,n})/\tilde n$.  Power conservation at a lossy
interface is the polarisation-wise identity
\begin{equation}\label{eq:fresnel-energy}
  T_{s,n}+|r_{s,n}|^2=1,\qquad T_{p,n}+|r_{p,n}|^2=1,
\end{equation}
where $T_{s,n}$ and $T_{p,n}$ are the monograph's power
transmittances.  This is exact for lossy half-spaces once we read
$T$ through the flux-weighted Poynting balance; see the appendix of
the appendix of the monograph that derives Fresnel coefficients
from the boundary conditions, and the discussion around the
power-transmittance definition $T=1-|r|^2$ in Part~I.

\subsection{A Fresnel reflection operator}

Mirror the construction of the Fresnel transmission operator
$\Fn(\rr)=\bigl(t_{s,n}\hat e_{s,n}\hat e_{s,n}^{\!T}
+t_{p,n}\hat e_{p,n}\hat e_{p,n}^{\!T}\bigr)H(\mu_n)$:

\begin{definition}[Fresnel reflection operator]\label{def:Rop}
At a front-facing surface point $\rr$ the per-path Fresnel
reflection operator is
\begin{equation}\label{eq:Rn}
  \Rn(\rr)=\Bigl(r_{s,n}\,\hat e_{s,n}\hat e_{s,n}^{\!T}
                +r_{p,n}\,\hat e_{p,n}\hat e_{p,n}^{\!T}\Bigr)H(\mu_n).
\end{equation}
\end{definition}

Just as $\Fn$ filters the incident polarisation $\bm\psi_n$ into the
transmitted surface field, $\Rn$ filters it into the specularly
reflected amplitude.  Both are rank-two, dimensionless, and differ
only in which Fresnel coefficient multiplies each polarisation
eigenvector.

\subsection{A scattered-power channel matrix and its Gram}

The monograph's exposure channel vector is
\begin{equation}
  \gtil_j(\rr)=\sum_{n:j(n)=j}\sqrt{\tfrac{\sigma}{4\alpha_n}}\,
  \Fn(\rr)\,\bm\psi_n\,e^{-ik_0\khat_n\cdot\rr}.
\end{equation}
The prefactor $\sqrt{\sigma/(4\alpha_n)}$ plus
the depth-integral identity of the monograph exactly turns
$|\Fn\bm\psi_n|^2$ into
$T_n\cos\theta_n/Z_0$.  For reflection we do not integrate over
depth, we evaluate the outgoing Poynting vector on $\Sigma^+$; the
amplitude-to-power conversion is just $1/Z_0$ times $\cos\theta_n$.
Define the \emph{scattered channel vector}
\begin{equation}\label{eq:gref}
  \gtil^{\mathrm{re}}_j(\rr)=\sum_{n:j(n)=j}
  \sqrt{\tfrac{\cos\theta_n}{Z_0}}\,\Rn(\rr)\,\bm\psi_n\,
  e^{-ik_0\khat_n^{\mathrm{re}}\cdot\rr},
\end{equation}
where $\khat_n^{\mathrm{re}}=\khat_n+2\mu_n\nhat$ is the specular
direction.  Stack into
$\Gtil^{\mathrm{re}}(\rr)=[\gtil^{\mathrm{re}}_1,\ldots,
\gtil^{\mathrm{re}}_M]\in\Complex^{3\times M}$.

\begin{definition}[Reflection operator]\label{def:Qref}
\begin{equation}\label{eq:Qref}
  \boxed{\Qref=\int_\Sigma\Gtil^{\mathrm{re}}(\rr)^H\,
                             \Gtil^{\mathrm{re}}(\rr)\,\diff A
         \;\in\;\Complex^{M\times M}.}
\end{equation}
\end{definition}

$\Qref$ is Hermitian positive-semidefinite, same proof as $\Qabs$.
Interpretation: $\mathbf{x}^H\Qref\mathbf{x}$ is the total power
scattered by the body surface into its exterior half-space, under
the same plane-wave superposition that defines $P_{\mathrm{ab}}$.

\subsection{The per-surface energy identity}

At a single front-facing point, Fresnel conservation
\cref{eq:fresnel-energy} is a \emph{per-path, per-polarisation}
statement.  When summed coherently over paths, the cross-terms
survive in both $\Sab$ and $\Sref$ in the same way:

\begin{proposition}[Local energy balance]\label{prop:local}
Under Approximations~1 and~2 of the monograph,
\begin{equation}\label{eq:local-balance}
  \Sab(\rr)+\Sref(\rr)=\Sinn(\rr),
\end{equation}
where $\Sref(\rr)=\|\Gtil^{\mathrm{re}}(\rr)\mathbf{x}\|^2$ is the
specularly scattered power density and
\begin{equation}\label{eq:Sinn}
  \Sinn(\rr)=\sum_n \frac{|x_{j(n)}|^2\cos\theta_n}{Z_0}
             |\bm\psi_n|^2 H(\mu_n)
             +\text{coherent cross-terms},
\end{equation}
is the \emph{incident} surface flux, i.e.\ the Poynting flux of
$\sum_n\EE_n$ through $\Sigma$.
\end{proposition}

Sketch.  The incident tangential and normal fields decompose into
TE and TM.  For TE, $|\psi_{s,n}|^2\cos\theta_n$ partitions as
$T_{s,n}|\psi_{s,n}|^2\cos\theta_n+|r_{s,n}|^2|\psi_{s,n}|^2\cos\theta_n$
by \cref{eq:fresnel-energy}, with identical angular and phase
factors on the two outgoing terms, so cross-terms obey the same
identity path-by-path.  TM is the same up to a factor of
$|\tilde n|^{-2}$ in the normal-field contribution, already absorbed
into the $O(|\tilde n|^{-2})$ error of Approximation~1.

Integrating \cref{eq:local-balance} over the body gives the
operator identity
\begin{empheq}[box=\tcbhighmath]{equation}\label{eq:op-balance}
  \Qinn=\Qabs+\Qref,\qquad
  \Qinn\equiv\int_\Sigma\mathbf{G}(\rr)^H\mathbf{G}(\rr)H(\mu)\,\diff A,
\end{empheq}
where $\mathbf{G}$ is the incident field channel matrix of the
monograph's field channel matrix.  $\Qinn$ is the
\emph{incidence operator}: it maps precoders to power that lands on
the body, irrespective of whether the body then absorbs it or
reflects it away.  $\Qabs$ and $\Qref$ partition that landed power.

\begin{remark}[What we have ignored]\label{rem:ignored}
Three contributions do not appear in \cref{eq:op-balance}.
(i) Evanescent/Norton tangential surface waves along the air-skin
interface; negligible at mmWave, order a few percent below~6\,GHz.
(ii) Creeping and diffraction around body contours, treated in
Part~I as the diffraction-smoothing correction.
(iii) Re-absorption after reflection, i.e.\ a path that reflects off
the torso, hits the arm, and is then absorbed there.  This is a
second-order term in the body's own albedo
$\bar\alpha=\tr\Qref/\tr\Qinn$ and is left for \cref{sec:higher}.
\end{remark}

% ============================================================
\section{Global book-keeping}
% ============================================================

\subsection{The missed operator}

What never hits the body?  Everything the transmitter radiates
outside the cone subtended by $\Sigma$.  Let $\mathbf{G}^{\mathrm{far}}$
be the channel matrix of the per-element far field evaluated on a
sphere at infinity; its Gram integrated over that sphere is a
multiple of~$\mathbf{I}$, because by reciprocity/unitarity the
total power radiated is $\|\mathbf{x}\|^2$ regardless of precoding.
Split the sphere into the solid angle that contains the body,
$\Omega_\Sigma$, and its complement.  Define
\begin{equation}\label{eq:Qmis}
  \Qmis=\|\mathbf{x}\|^{-2}\!\!\int_{\Omega_\Sigma^c}\!\!
        \mathbf{G}^{\mathrm{far}}(\Omega)^H\mathbf{G}^{\mathrm{far}}(\Omega)\,\diff\Omega.
\end{equation}
This operator is cleanly defined when the body is the only scatterer.
With walls and buildings, $\Qmis$ absorbs into a bigger
\emph{environmental loss} operator $\mathbf{Q}_{\mathrm{env}}$
(dissipation in concrete and glass) plus the true sky loss.  We
keep the simple form.

\subsection{The bound}

In free space with a single scatterer, reciprocity gives

\begin{theorem}[Array-level power conservation]\label{thm:global}
In a lossless background with the body as the only scatterer, for
every precoder $\mathbf{x}\in\Complex^M$,
\begin{equation}\label{eq:global-balance}
  \mathbf{x}^H\bigl(\Qabs+\Qref+\Qmis\bigr)\mathbf{x}
  =\|\mathbf{x}\|^2,
\end{equation}
equivalently $\Qabs+\Qref+\Qmis=\mathbf{I}$ as Hermitian forms.
With additional loss channels (walls, etc.), $\preceq$ replaces $=$.
\end{theorem}

Proof sketch.  The total radiated power is $\|\mathbf{x}\|^2$ and
flows through any enclosing surface.  Split the surface into the
body and a sphere at infinity.  On the body, Poynting in equals
$\mathbf{x}^H\Qinn\mathbf{x}=\mathbf{x}^H(\Qabs+\Qref)\mathbf{x}$ by
\cref{eq:op-balance}, where ``in'' is net inflow counting the
reflected beam as part of the outflow on the far sphere.  The rest,
which reaches infinity directly without interacting with the body,
is $\mathbf{x}^H\Qmis\mathbf{x}$.  By conservation these three sum
to $\|\mathbf{x}\|^2$.

\Cref{thm:global} is the reason the title of this note calls
$\Qref+\Qmis$ \emph{the complement}: it is literally
$\mathbf{I}-\Qabs$, modulo environmental loss.

\begin{corollary}[Kirchhoff-like identity for modes]\label{cor:modes}
Let $\lambda_k^{\mathrm{ab}}$, $\lambda_k^{\mathrm{re}}$,
$\lambda_k^{\mathrm{mi}}$ be the eigenvalues of $\Qabs,\Qref,\Qmis$
along a common precoder direction $\mathbf{v}_k$.  Then
$\lambda_k^{\mathrm{ab}}+\lambda_k^{\mathrm{re}}
+\lambda_k^{\mathrm{mi}}=1$.
\end{corollary}

This is a literal sum-of-probabilities over ``absorbed / reflected /
missed'' if we think of precoders as pure quantum-mechanical states
and the three operators as the POVM of outcomes.  It is almost too
tidy.

% ============================================================
\section{Pseudo-Brewster eigenvector locking}
% ============================================================

Now the fun starts.  $\Qabs$ and $\Qref$ are built from the same
paths, the same body, the same precoder, with different scalar
weights ($t_s,t_p$ vs $r_s,r_p$) on the same TE/TM polarisation
bases.  What if those weights are nearly the same up to an overall
constant?

\subsection{The observation}

Under pseudo-Brewster (Part~I of the monograph) we have
$T_s(\theta)\approx T_p(\theta)\approx T_0$ across
$\theta\in[0,80^\circ]$ for skin and other high-water tissue.  Hence
$|r_s|^2\approx|r_p|^2\approx 1-T_0$, and more strongly, the
amplitude coefficients are nearly equal:
\begin{equation}\label{eq:pB-amp}
  t_{s,n}\approx t_{p,n}\approx\sqrt{T_0}\,e^{i\phi_t(\theta_n)},
  \qquad
  r_{s,n}\approx r_{p,n}\approx\sqrt{1-T_0}\,e^{i\phi_r(\theta_n)}.
\end{equation}
The phases $\phi_t,\phi_r$ are angle-dependent but
polarisation-blind in this regime.  Substituting into \cref{eq:Rn}
and the analogous $\Fn$,
\begin{equation}\label{eq:F-R-collapse}
  \Fn\approx \sqrt{T_0}\,e^{i\phi_t}\Pi_n,\qquad
  \Rn\approx \sqrt{1-T_0}\,e^{i\phi_r}\Pi_n,
\end{equation}
where $\Pi_n=\hat e_{s,n}\hat e_{s,n}^{\!T}+\hat e_{p,n}\hat e_{p,n}^{\!T}
=H(\mu_n)(\mathbf{I}_3-\khat_n\khat_n^{\!T})$ is the polarisation
projector transverse to $\khat_n$.  That is the key identity.

\subsection{Consequence: shared eigenvectors}

\begin{proposition}[Eigenvector locking]\label{prop:lock}
Under the pseudo-Brewster amplitude approximation \cref{eq:pB-amp},
\begin{equation}\label{eq:lock}
  \Qref\approx\frac{1-T_0}{T_0}\,\Qabs + \Delta,
\end{equation}
where $\|\Delta\|/\|\Qabs\|=O(\max|T_s-T_p|/T_0)\lesssim 5\%$ at mmWave.
In particular, $\Qabs$ and $\Qref$ share their top eigenvectors
(to within that tolerance), and for every unit precoder
$\mathbf{x}$,
\begin{equation}\label{eq:lock-ratio}
  \mathbf{x}^H\Qref\mathbf{x}\approx
  \frac{1-T_0}{T_0}\,\mathbf{x}^H\Qabs\mathbf{x}.
\end{equation}
\end{proposition}

Proof sketch.  In \cref{eq:F-R-collapse} the phase factors
$e^{i\phi_t},e^{i\phi_r}$ depend on $\theta_n$ but not on
polarisation; they multiply each per-path contribution to
$\Gtil(\rr)$ and $\Gtil^{\mathrm{re}}(\rr)$ by the \emph{same} scalar
per path.  The two Gram integrals then differ only by the replacement
$T_0\leftarrow 1-T_0$ and by the outgoing direction $\khat_n$ vs
$\khat_n^{\mathrm{re}}$.  The outgoing-direction change affects only
the per-point phase $e^{-ik_0\khat\cdot\rr}$, which is a unitary
factor on the channel matrix: it leaves
$\mathbf{G}^{\!H}\mathbf{G}$ invariant up to the sign of
cross-element phases that integrate out on $\Sigma$.  What remains
is the scalar factor $(1-T_0)/T_0$.

Equation \cref{eq:lock-ratio} is striking.  \emph{The worst-case
precoder for body exposure is the worst-case precoder for body
scattering.}  In the pseudo-Brewster regime these two optimisation
problems are the same problem with a rescaled objective.

\subsection{Why this is a nontrivial statement}

Pseudo-Brewster already told us that $\Sab\approx\Sinc\,T_0\,\ReLU(\mu)$
is polarisation-blind per-wave.  What's new here is at the
\emph{array level}: the entire Hermitian operator $\Qabs$ is
proportional to $\Qref$, with a body-universal proportionality
constant $(1-T_0)/T_0\approx 0.63/0.37\approx 1.7$ for skin at
28\,GHz.  So every exposure mode $\mathbf{v}_k$ of $\Qabs$ is also a
scattering mode, with scattered power
$\lambda_k^{\mathrm{re}}\approx 1.7\lambda_k^{\mathrm{ab}}$.

Two paragraphs ago this was a bookkeeping exercise.  Now it has a
punch line: the body cannot be dosed more without being lit up more
on a monostatic radar.

% ============================================================
\section{Sensing-communication duality, sketched}\label{sec:jsac}
% ============================================================

This is where the complement matrix earns its keep.  In joint
sensing and communication (JSAC), we transmit $\mathbf{x}$ to serve
a UE with channel $\mathbf{h}$ \emph{and} to image the environment
via backscatter; the body is the dominant near-array scatterer.
Sensing gain for a colocated mono-static receiver is
\begin{equation}\label{eq:jsac-sense}
  G_{\mathrm{sense}}(\mathbf{x})=
  \mathbf{x}^H\Qref\mathbf{x}\cdot\mathcal{K}(\mathbf{h}_{\mathrm{rx}}),
\end{equation}
with $\mathcal{K}$ the receive-aperture collection factor we do not
need to unpack here.  Exposure-constrained beamforming (Part~III of the monograph) solves
\begin{equation}\label{eq:ecbf-orig}
  \max_{\mathbf{x}}\;|\mathbf{h}^H\mathbf{x}|^2
  \quad\text{s.t.}\quad \mathbf{x}^H\Qabs\mathbf{x}\le P_0.
\end{equation}
Under \cref{prop:lock} the feasibility set
$\{\mathbf{x}:\mathbf{x}^H\Qabs\mathbf{x}\le P_0\}$ is, up to
$5\,\%$, the same as
$\{\mathbf{x}:\mathbf{x}^H\Qref\mathbf{x}\le P_0\cdot(1-T_0)/T_0\}$.
Two consequences:

\begin{enumerate}[leftmargin=1.6em,itemsep=2pt]
\item \textbf{Safe-JSAC is a bound constraint, not a trade-off.}
      Shrinking exposure shrinks sensing return proportionally
      because the two Gram forms are parallel.  You do not pay for
      ``communication + sensing'' against ``exposure'' as three
      axes; exposure and sensing live on the same axis.
\item \textbf{Nullspace precoders are the escape.}  The only
      $\mathbf{x}$ that deliver $\mathbf{h}^H\mathbf{x}\ne 0$ with
      $\mathbf{x}^H\Qabs\mathbf{x}=0$ are those with
      $\mathbf{x}\in\ker\Qinn$ (the transparency subspace below).
      These carry information to the UE without touching the body
      and without scattering off it either.  Perfect for ``quiet''
      communication, useless for imaging.
\end{enumerate}

I do not yet know how large $\ker\Qinn$ is for typical array sizes
and body poses.  It is clearly non-empty when $M$ exceeds the number
of independent polarised paths through $\Sigma$, which for a body
and 256-element arrays is the common case.  The dimensionality of
$\ker\Qinn$ is the key quantity for quiet JSAC and I mark it for
follow-up.

% ============================================================
\section{The transparency subspace}
% ============================================================

\subsection{Definition and dimension}

\begin{definition}[Transparency subspace]
\begin{equation}
  \mathcal{T}\equiv\ker\Qinn=
  \{\mathbf{x}\in\Complex^M:\Qinn\mathbf{x}=\mathbf{0}\}.
\end{equation}
\end{definition}

Because $\Qinn$ is a Gram, $\mathcal{T}$ is the orthogonal
complement of the row space of the stacked channel matrix
$[\mathbf{G}(\rr_1);\ldots;\mathbf{G}(\rr_K)]$ sampled on a
sufficiently dense quadrature of $\Sigma$.  Its dimension is
$M-\mathrm{rank}\,\Qinn$.  A precoder $\mathbf{x}\in\mathcal{T}$
produces zero incident surface field on $\Sigma$, therefore zero
absorbed and zero scattered power.  From the body's point of view
the array is silent.

\begin{proposition}[Generic transparency]\label{prop:dim}
Let $N_{\mathrm{body}}$ be the number of distinct propagation paths
reaching any surface element of $\Sigma$ from at least one array
element (typically $N_{\mathrm{body}}\le 3M_\triangle$ where
$M_\triangle$ is the number of front-facing triangles).  Then
\begin{equation}\label{eq:dim}
  \dim\mathcal{T}\ge\max(0,\,M-N_{\mathrm{body}}).
\end{equation}
\end{proposition}

The bound is typically loose.  Inter-element path correlation, the
finite solid angle subtended by the body, and the three-dimensional
polarisation structure drive $\mathrm{rank}\,\Qinn$ well below its
nominal upper bound.  A numerical measurement of $\dim\mathcal{T}$
for, say, a 256-element URA at 28\,GHz facing a thelonious phantom
at 1\,m is a one-afternoon experiment and I expect
$\dim\mathcal{T}/M\gtrsim 0.3$ to $0.7$ depending on pose.  Flag:
this is a conjecture based on how the monograph's directivity spectra
decay, not a measurement.

\subsection{The soft version}

In practice we do not want a precoder in exact $\ker\Qinn$, we want
one in the \emph{low-absorption eigenspan}.  The matrix
\begin{equation}
  \mathbf{P}_\varepsilon =
  \sum_{k:\lambda_k^{\mathrm{ab}}/\lambda_1^{\mathrm{ab}}\le\varepsilon}
        \mathbf{v}_k\mathbf{v}_k^{\!H}
\end{equation}
projects onto the $\varepsilon$-soft transparency subspace.  The
monograph's ECBF solution is obtained by steering
$(\lambda\Qabs+\nu\mathbf{I})^{-1}\mathbf{h}^*$; the transparency
view is the complementary picture, where we first kill everything
outside $\mathbf{P}_\varepsilon$ and then find the best
UE-alignment inside it.  The two routes give the same optimiser in
the limits, but the transparency picture makes the geometry of the
feasibility set explicit.

% ============================================================
\section{Kirchhoff dual: $\Qabs$ is also a receiver}
% ============================================================

The same operator $\Qabs$ that counts absorbed power on a transmit
precoder $\mathbf{x}$ counts, by reciprocity, the coupling of
thermal emission from the body into a \emph{receive} combiner.
This is Kirchhoff in array form.

\subsection{Absorptivity equals emissivity, operator version}

For a body at temperature $T$, spectral radiance per unit area per
solid angle per polarisation is $B_\lambda(T)\cdot\epsilon(\rr,\khat)$
with $\epsilon=1-|r|^2$.  By \cref{eq:fresnel-energy},
$\epsilon=T_{s,p}$ per polarisation.  So the per-surface element
emissive power toward direction $\khat'$ factors through the same
$\Fn$-like operator that defined absorption, evaluated on the
\emph{outgoing} path $(\rr,\khat')$.

\begin{proposition}[Reciprocal coupling]\label{prop:kirch}
Let $\mathbf{y}\in\Complex^M$ be a receive combiner applied to the
array with effective noise temperature dominated by body radiation.
Up to a Planck prefactor $\Phi(f,T)$,
\begin{equation}\label{eq:kirch}
  \langle|\mathbf{y}^H\mathbf{n}_{\mathrm{body}}|^2\rangle
  \approx \Phi(f,T)\cdot \mathbf{y}^H\Qabs\mathbf{y}.
\end{equation}
Thermal pick-up is governed by the \emph{transmit} exposure operator.
\end{proposition}

Sketch.  Reciprocity for the two-port system ``array'' and
``each body surface patch'' maps transmitted amplitude through $\Fn$
to received amplitude through $\Fn^{\!T}$; the Gram integral over
$\Sigma$ is the same in both directions.  The Planck prefactor
collects bandwidth and temperature.

\subsection{What this means in practice}

If an array is precoded to maximise body exposure (top eigenvector
of $\Qabs$), the conjugate combiner is the matched filter for
body thermal emission.  Passive millimetre-wave body imaging and
active mmWave exposure share the same spectral basis.  Exposure
control and thermographic sensitivity are reciprocal problems.

% ============================================================
\section{A trace-level compliance law}
% ============================================================

The monograph gives $\lambda_1(\Qabs)$ as the worst-case
whole-body absorbed power per unit transmit power, following
(Ebadi-Shahrivar \textit{et al.}, 2019).  The complement view adds
three free identities:

\begin{description}[leftmargin=0pt,style=nextline]
\item[Albedo.] The body's mean albedo under isotropic transmit is
\[
  \bar\alpha = \frac{\tr\Qref}{\tr\Qinn}
  \approx 1-T_0\approx 0.37 \text{ at mmWave}.
\]
\item[Capture cross-section.] The body's effective capture
fraction of the array's radiated power is
\[
  \bar\kappa = \frac{\tr\Qinn}{M}.
\]
This is the scaled isotropic-transmit solid-angle coverage, modulated
by pseudo-Brewster.  Limit $\bar\kappa\to\Omega_\Sigma/4\pi$ in the
incoherent regime.
\item[Array-level compliance.]  ICNIRP-style whole-body SAR
compliance is $\lambda_1(\Qabs)\le P_0\cdot m_{\mathrm{body}}\cdot
\mathrm{SAR}_{\mathrm{wb,lim}}/P_{\mathrm{tx}}$.  Under
\cref{prop:lock}, this is equivalent to a bound on
$\lambda_1(\Qref)$ with the body-universal constant $(1-T_0)/T_0$.
Two constraints, one eigenvalue.
\end{description}

% ============================================================
\section{Where the book-keeping breaks down}\label{sec:higher}
% ============================================================

This whole note lives in the single-bounce pseudo-Brewster regime.
Three obvious directions in which the clean picture loosens:

\begin{enumerate}[leftmargin=1.6em,itemsep=2pt]
\item \textbf{Multi-bounce body reflections.}  The torso scatters
onto the arm; the arm then absorbs some of that.  A Neumann series
on $\Qref$ gives the correction:
\[
  \Qabs^{(\infty)} = \Qabs + \Qabs\cdot\mathcal{R}\cdot\Qref
  + \Qabs\cdot\mathcal{R}\cdot\Qref\cdot\mathcal{R}\cdot\Qref+\cdots,
\]
where $\mathcal{R}$ is a body-to-body visibility operator.  This is
the classical radiosity cascade of Part~I, lifted to the coherent
MIMO level.  Convergence is guaranteed by $\|\mathcal{R}\Qref\|\le
\bar\alpha\cdot V_{\mathrm{self}}$ with self-visibility
$V_{\mathrm{self}}$ of the phantom, typically $\ll 1$.
\item \textbf{Below-6-GHz regime.}  Pseudo-Brewster is accurate only
at mmWave.  At sub-6, \cref{prop:lock} becomes an inequality
$\Qref\approx c_s(\theta)\Qabs^{(s)}+c_p(\theta)\Qabs^{(p)}$ with
polarisation-dependent coefficients; the shared-eigenvector result
becomes a block-sharing result.
\item \textbf{Non-planar skin.}  Curvature and wavelength-comparable
features give diffraction, not Fresnel.  $\Qabs+\Qref\ne\Qinn$
there; the defect is the monograph's diffraction-smoothing term.
\end{enumerate}

% ============================================================
\section{Loose threads}
% ============================================================

A list of leads the bookkeeping has opened.  Not claims.  Questions.

\begin{enumerate}[leftmargin=1.6em,itemsep=3pt]
\item Does $\Qref$ admit a closed-form bound analogous to
$\lambda_1(\Qabs)\le T_0\cdot\Omega_{\mathrm{body}}/Z_0$ from
Part~I?  \cref{prop:lock} suggests
$\lambda_1(\Qref)\le(1-T_0)\cdot\Omega_{\mathrm{body}}/Z_0$
with the same geometric factor.  Check.
\item Is there a clean ECBF dual
$\max_{\mathbf{x}}\mathbf{x}^H\Qref\mathbf{x}
\text{ s.t. }\mathbf{x}^H\Qabs\mathbf{x}\le P_0$?  Under
\cref{prop:lock} it degenerates to the trivial
$\|\mathbf{x}\|^2\le P_0 T_0/(1-T_0)$, but away from pseudo-Brewster
it is a real problem with its own optimiser, and it describes the
best body-illuminator under a fixed dose budget.  Medical imaging
reads this as ``best contrast per absorbed dose''.
\item The transparency subspace dimension $\dim\mathcal{T}$ for
representative base-station poses.  I would like a figure:
cumulative eigenvalue fraction of $\Qinn$ vs mode index, overlaid
for 10, 100, 1000-element arrays.  Conjecture: rapid decay; the
``exposure information'' of a body is low-dimensional.
\item Multi-user MU-MIMO.  Each user $u$ has a per-person operator
$\mathbf{Q}^{(u)}$ (Part~III, multi-user section).  Do the
per-person reflection operators $\mathbf{Q}^{(u)}_{\mathrm{re}}$
still interleave with their absorption counterparts?  The
pseudo-Brewster argument is per-material and per-body, so yes, but
the cross-person $\mathbf{Q}^{(u,v)}_{\mathrm{re}}$ term (user $v$'s
reflection hitting user $u$) is a new object.
\item The thermal-dual statement \cref{eq:kirch} turns $\Qabs$ into
a passive-sensing operator.  Is there a calibration procedure that
\emph{measures} $\Qabs$ in situ by listening to body thermal
emission through the array?  That would turn compliance into
something an array could self-verify, without phantom.
\item The JSAC duality.  Given that exposure and backscatter share
eigenvectors, is there a regulator-level argument for
\emph{simultaneous} compliance on $\lambda_1(\Qabs)$ and
intentional sensing?  The radar engineer wants large
$\lambda_1(\Qref)$; the dosimetrist wants small $\lambda_1(\Qabs)$;
the pseudo-Brewster identity says you can't have one without the
other, so regulation can be stated on either.
\end{enumerate}

% ============================================================
\section{Summary}
% ============================================================

Start from the monograph's exposure operator $\Qabs$ and ask where
the un-absorbed watts live.  Three Hermitian companions fall out:
$\Qref$ (specular scattering), $\Qmis$ (radiation that never hit the
body), and their joint book-keeping identity $\Qabs+\Qref+\Qmis=
\mathbf{I}$.  Fresnel energy conservation localises on
$\Qinn=\Qabs+\Qref$.  Pseudo-Brewster then performs the interesting
collapse: $\Qref\approx(1-T_0)/T_0\cdot\Qabs$, shared eigenvectors,
same modes.  The complement of the exposure operator is not a new
degree of freedom.  It is the same geometry, seen from the other
side of the skin.

\end{document}
